% c00.fig.hohlraum — the usual realisation of a black body: a box with a small hole.
%
% A scheme.  Content: lecture notes vl01.tex:47 ("Kiste mit kleinem Loch: Einfallende
% Strahlung verschwindet innen durch Vielfachreflexion").
% Geometry: a box [0,4] x [0,2.4] (cm) with a hole in the left wall between y = 1.1 and
% y = 1.3.  One ray enters through the middle of the hole, (0, 1.2), with direction
% (1, 0.62) and is reflected specularly (angle of incidence = angle of reflection: at a
% horizontal wall the y-component of the direction changes sign, at a vertical wall the
% x-component).  The successive points of impact computed from this rule are
%   (1.9355, 2.4), (4, 1.12), (2.1935, 0), (0, 1.36), (1.6774, 2.4), (4, 0.96), (2.4516, 0):
% The drawn reflections do not hit the hole. At the drawn left-wall impact y = 1.36,
% the ray misses the hole; continuation can return to it. The last segment is dotted.
% The drawing only shows the ray path; the physical cavity idealisation needs an
% author-supplied source or wording (G-0024).
\begin{tikzpicture}[x=1cm, y=1cm]
  \fill[tpc bath] (-0.18,-0.18) rectangle (4.18,2.58);
  \fill[tpc-plate] (0,0) rectangle (4,2.4);
  \fill[tpc-plate] (-0.2,1.1) rectangle (0.01,1.3);
  \draw[tpc bath edge] (-0.18,1.1) -- (-0.18,-0.18) -- (4.18,-0.18) -- (4.18,2.58) -- (-0.18,2.58) -- (-0.18,1.3);
  \draw[tpc axis] (0,1.1) -- (0,0) -- (4,0) -- (4,2.4) -- (0,2.4) -- (0,1.3);
  \draw[tpc dynamics, -{Stealth[length=5pt]}] (-1.6,0.208) -- (-0.4,0.952);
  \draw[tpc dynamics] (-0.4,0.952) -- (0,1.2) -- (1.9355,2.4) -- (4,1.12) -- (2.1935,0) -- (0,1.36)
    -- (1.6774,2.4) -- (4,0.96);
  \draw[tpc dynamics, dotted] (4,0.96) -- (2.4516,0);
  \node[tpc label, anchor=east] at (-1.65,0.2) {\tpcifenglish{incident radiation}{einfallende Strahlung}};
  \draw[tpc leader] (-0.1,1.34) -- (-0.7,1.9);
  \node[tpc label, anchor=east] at (-0.72,1.95) {\tpcifenglish{hole}{Loch}};
  \node[tpc label, anchor=west] at (4.35,1.2) {\tpcifenglish{walls at temperature}{Wände mit Temperatur} $T$};
\end{tikzpicture}
